% ================================================================
% Portada del libro -- 100% TikZ, pagina completa
% ================================================================
\begin{titlepage}
\thispagestyle{empty}
\begin{tikzpicture}[remember picture, overlay]

  % ------------------------------------------------------------
  % Fondo: malla triangular tipo discretizacion FEM con degradado
  % (simula un mapa de contornos de tension), tercio inferior.
  % ------------------------------------------------------------
  \begin{scope}[shift={(current page.south west)}]
    \pgfmathsetmacro{\cellw}{2.45}
    \pgfmathsetmacro{\cellh}{3.15}
    \foreach \i in {0,...,8} {
      \foreach \j in {0,...,2} {
        \pgfmathsetmacro{\ax}{\i*\cellw + 0.14*sin(137*\i+91*\j)}
        \pgfmathsetmacro{\ay}{\j*\cellh + 0.14*cos(83*\i+59*\j)}
        \pgfmathsetmacro{\bx}{(\i+1)*\cellw + 0.14*sin(137*(\i+1)+91*\j)}
        \pgfmathsetmacro{\by}{\j*\cellh + 0.14*cos(83*(\i+1)+59*\j)}
        \pgfmathsetmacro{\cx}{\i*\cellw + 0.14*sin(137*\i+91*(\j+1))}
        \pgfmathsetmacro{\cy}{(\j+1)*\cellh + 0.14*cos(83*\i+59*(\j+1))}
        \pgfmathsetmacro{\dx}{(\i+1)*\cellw + 0.14*sin(137*(\i+1)+91*(\j+1))}
        \pgfmathsetmacro{\dy}{(\j+1)*\cellh + 0.14*cos(83*(\i+1)+59*(\j+1))}
        \pgfmathsetmacro{\mixpct}{100*\j/2}
        \fill[ucred!\mixpct!ucblue, opacity=0.55] (\ax,\ay) -- (\bx,\by) -- (\dx,\dy) -- cycle;
        \fill[ucred!\mixpct!ucblue, opacity=0.55] (\ax,\ay) -- (\dx,\dy) -- (\cx,\cy) -- cycle;
        \draw[white, opacity=0.55, line width=0.25pt] (\ax,\ay) -- (\bx,\by) -- (\dx,\dy) -- (\cx,\cy) -- cycle;
        \draw[white, opacity=0.55, line width=0.25pt] (\ax,\ay) -- (\dx,\dy);
      }
    }
    % Velo blanco tenue para que el texto inferior sea legible sobre la malla
    \fill[white, opacity=0.12] (0,0) rectangle (22.5,9.7);
  \end{scope}

  % ------------------------------------------------------------
  % Cubo de tensiones 3D (motivo "hero"): medio continuo -> FEM
  % Proyeccion isometrica via remapeo de ejes x/y/z (sin libreria 3d)
  % ------------------------------------------------------------
  \begin{scope}[shift={($(current page.north west)+(6.4,-8.9)$)}]
    \begin{scope}[x={(-0.866cm,-0.5cm)}, y={(0.866cm,-0.5cm)}, z={(0cm,1cm)}]
      \pgfmathsetmacro{\Lc}{2.3}
      \draw[fill=ucblue!55, draw=ucblue!80!black, line width=0.6pt]
        (0,0,\Lc) -- (\Lc,0,\Lc) -- (\Lc,\Lc,\Lc) -- (0,\Lc,\Lc) -- cycle; % cara superior
      \draw[fill=ucblue!75, draw=ucblue!80!black, line width=0.6pt]
        (\Lc,0,0) -- (\Lc,\Lc,0) -- (\Lc,\Lc,\Lc) -- (\Lc,0,\Lc) -- cycle; % cara derecha
      \draw[fill=ucblue!65, draw=ucblue!80!black, line width=0.6pt]
        (0,\Lc,0) -- (\Lc,\Lc,0) -- (\Lc,\Lc,\Lc) -- (0,\Lc,\Lc) -- cycle; % cara frontal
      % Vectores de traccion normales a cada cara visible
      \draw[-{Stealth[length=2.6mm]}, ucred, line width=0.9pt]
        (\Lc/2,\Lc/2,\Lc) -- ++(0,0,0.85) node[above, ucred, font=\scriptsize] {$\bm{t}$};
      \draw[-{Stealth[length=2.6mm]}, ucred, line width=0.9pt]
        (\Lc,\Lc/2,\Lc/2) -- ++(0.85,0,0) node[right, ucred, font=\scriptsize] {$\bm{t}$};
      \draw[-{Stealth[length=2.6mm]}, ucred, line width=0.9pt]
        (\Lc/2,\Lc,\Lc/2) -- ++(0,0.85,0) node[above right, ucred, font=\scriptsize] {$\bm{t}$};
    \end{scope}
  \end{scope}

  % ------------------------------------------------------------
  % Encabezado institucional
  % ------------------------------------------------------------
  \node[anchor=north, align=center, text width=0.8\paperwidth] at ($(current page.north)+(0,-1.9)$) {%
    \scshape\large Pontificia Universidad Católica de Chile\\[2pt]
    \scshape\normalsize Escuela de Ingeniería\\[2pt]
    \upshape\small Departamento de Ingeniería Mecánica y Metalúrgica\\[2pt]
    \upshape\small ICM 3323 --- Elementos Finitos en Diseño Mecánico
  };
  \draw[ucred, line width=0.6pt]
    ($(current page.north)+(-0.21\paperwidth,-3.05)$) -- ($(current page.north)+(0.21\paperwidth,-3.05)$);

  % Insignia generica (no es el escudo oficial de la universidad)
  \node[circle, draw=ucred, line width=0.8pt, minimum size=1.35cm, fill=white]
    at ($(current page.north)+(0,-1.05)$) {\scshape\footnotesize\color{ucred} PUC};

  % ------------------------------------------------------------
  % Titulo y subtitulo
  % ------------------------------------------------------------
  \node[anchor=north, align=center, text width=0.82\paperwidth] at ($(current page.north)+(0,-13.6)$) {%
    {\fontsize{30}{34}\selectfont\bfseries\color{ucred} Elementos Finitos\\[2pt] en Diseño Mecánico}
  };
  \node[anchor=north, align=center, text width=0.72\paperwidth] at ($(current page.north)+(0,-16.1)$) {%
    {\Large\color{ucblue} Mecánica del Medio Continuo y Formulación\\ del Método de los Elementos Finitos}\\[4pt]
    {\normalsize\itshape\color{black!65} Apuntes de cátedra}
  };

  % ------------------------------------------------------------
  % Bloque inferior: autoria, curso, fecha
  % ------------------------------------------------------------
  \node[anchor=south, align=center, text width=0.78\paperwidth] at ($(current page.south)+(0,2.55)$) {%
    {\small Basado en los apuntes del curso dictado por el\\ Prof.\ Diego Celentano}\\[8pt]
    {\normalsize\bfseries Recopilado y editado por: Pedro Taboada}\\[3pt]
    {\small Primer Semestre 2026 --- Pontificia Universidad Católica de Chile}
  };
  \draw[ucred, line width=1.4pt] ($(current page.south)+(-0.30\paperwidth,1.65)$) -- ($(current page.south)+(0.30\paperwidth,1.65)$);

\end{tikzpicture}
\end{titlepage}
